\section{Vocabulario y catálogos}\label{sec:vocabulario}

\subsection{Conceptos y propiedades}
Escribo un concepto indexado con K01; representa una clase de objetos, cuyo
conjunto de miembros se determina en el contexto. Una propiedad indexada usa K02
y representa un criterio que puede cumplirse o no. Los índices se conservan al
cambiar sus etiquetas explicativas. Ninguna etiqueta introduce un sonido.

\begin{longtable}{@{}p{26mm}>{\raggedright\arraybackslash}p{\dimexpr\linewidth-26mm-2\tabcolsep\relax}@{}}
\toprule Escritura & Definición y dominio \\ \midrule\endhead
\KalaConcepto{1} & C1. Personas identificadas en el estudio declarado. \\[5pt]
\KalaConcepto{2} & C2. Libros identificados en ese estudio; se declara si se cuentan obras, ediciones o ejemplares. T04--T05 usan ejemplares. \\[5pt]
\KalaConcepto{3} & C3. Sensores del inventario técnico declarado. \\[5pt]
\KalaConcepto{4} & C4. Muestras materiales individualizadas en el registro experimental. \\[5pt]
\KalaConcepto{5} & C5. Tonos estables del intervalo de observación declarado. \\[5pt]
\KalaPropiedad{1} & P1. Rojo: cumple un criterio de clasificación de color especificado en el contexto, con condiciones de observación. \\[5pt]
\KalaPropiedad{2} & P2. Activo: el dispositivo tiene el estado operativo designado como activo en el instante de inspección declarado. \\[5pt]
\KalaPropiedad{3} & P3. Líquido: la muestra presenta la fase clasificada como líquida, bajo la temperatura y presión declaradas. \\[5pt]
\bottomrule
\end{longtable}
Las letras C y P de esta explicación son ayudas de catálogo. La escritura usa
el signo seguido del índice. La proposición convencional \(P\) de T02 no es
la propiedad P1: pertenecen a construcciones distintas.

\subsection{Colecciones}
K06 seguido del índice \textbf{1} y un grupo construye un conjunto finito por
enumeración. El orden no cambia el conjunto y repetir un miembro no lo añade dos
veces. El índice \textbf{2} construye una lista finita: conserva orden y repeticiones.
Un grupo vacío construye el conjunto vacío o la lista vacía según el índice.
Los elementos se separan con la misma barra vertical de argumentos.

\KalaMostrar{\KalaColeccion{1}{\KalaReferencia{1}\KalaSeparador\KalaReferencia{2}}
\qquad\KalaColeccion{2}{\KalaReferencia{1}\KalaSeparador\KalaReferencia{2}}}
En esta ilustración los referentes 1 y 2 son objetos distintos A y B.
Los signos K06 y K20 no son intercambiables: K06 enumera miembros; K20 selecciona
miembros de un dominio por una condición. Ambos pueden producir conjuntos.

\subsection{Acción elemental del fragmento de programas}
K04 seguido de \textbf{1} y dos argumentos indica asignación:
variable de destino y expresión entera. El lado derecho se evalúa en el estado
anterior; luego se actualiza solo la variable de destino. Su resultado es un
programa, no una proposición de igualdad. No modifica la interpretación de K04
sin índice, que sigue nombrando la categoría proceso.

\KalaMostrar{\KalaAsignar{x}{x+1}}
Ejemplo: si \(x=2\), la asignación deja \(x=3\). El índice de acción 1 es
independiente del índice de relación 1, del referente 1 y del concepto 1.

\subsection{Regla para ampliar un catálogo}
Registrar: identificador sin reutilizar, tipo de expresión, definición, dominio,
argumentos ordenados, contexto obligatorio, ejemplo positivo y un contraste.
Un nombre corto podrá abreviar una composición después de definirla, sin cambiar
su significado. Los catálogos se versionan junto con la gramática.
